\documentclass[12pt,a4paper]{article}
\usepackage[utf8]{inputenc}
\usepackage[T1]{fontenc}
\usepackage[spanish]{babel}
\usepackage{amsmath,amssymb}
\usepackage{geometry}
\usepackage{booktabs}
\usepackage{array}
\usepackage{xcolor}
\usepackage{tcolorbox}
\usepackage{fancyhdr}
\usepackage{lmodern}
\usepackage{microtype}
\usepackage{enumitem}
\usepackage{tikz}
\usetikzlibrary{arrows.meta, positioning, shapes.geometric}

\geometry{top=2.5cm, bottom=2.5cm, left=2.5cm, right=2.5cm, headheight=25pt}

\definecolor{azulul}{RGB}{0,70,127}
\definecolor{verdeclaro}{RGB}{200,230,201}
\definecolor{azulclaro}{RGB}{200,220,240}
\definecolor{amarilloclaro}{RGB}{255,249,196}

\tcbuselibrary{skins, breakable}

\pagestyle{fancy}
\fancyhf{}
\lhead{\color{azulul}\textbf{Econom\'ia Financiera 1 -- Gu\'ia 5}}
\rhead{\color{azulul}Ejercicio 2: Industrias LAN S.A.}
\cfoot{\thepage}
\renewcommand{\headrulewidth}{1pt}
\renewcommand{\headrule}{\hbox to\headwidth{\color{azulul}\leaders\hrule height \headrulewidth\hfill}}

\newtcolorbox{cajaenunciado}{
  colback=azulclaro!40, colframe=azulul, fonttitle=\bfseries,
  title=Enunciado del Ejercicio 2, breakable
}
\newtcolorbox{cajaintuicion}[1]{
  colback=verdeclaro!60, colframe=green!50!black, fonttitle=\bfseries,
  title={Intuici\'on: #1}, breakable
}
\newtcolorbox{cajaresultado}{
  colback=verdeclaro!80, colframe=green!40!black, fonttitle=\bfseries\large,
  title=Resultado Final, halign title=center
}
\newtcolorbox{cajaformula}{
  colback=azulclaro!20, colframe=azulul!80, fonttitle=\bfseries,
  title=F\'ormula clave, breakable
}
\newtcolorbox{cajareflexion}[1]{
  colback=amarilloclaro, colframe=orange!70!black, fonttitle=\bfseries,
  title={Reflexi\'on: #1}, breakable
}

\begin{document}

% ========================
%  PORTADA
% ========================
\begin{titlepage}
\centering
\vspace*{1cm}

{\color{azulul}\rule{\linewidth}{1.5pt}}
\vspace{0.5cm}

{\large\textbf{UNIVERSIDAD DE LIMA}\\[0.2cm]
Facultad de Ciencias Empresariales y Econ\'omicas\\
Carrera de Econom\'ia}

\vspace{1.5cm}

{\LARGE\color{azulul}\textbf{ECONOM\'IA FINANCIERA 1}}\\[0.4cm]
{\Large Gu\'ia de Pr\'actica 5 --- Pol\'itica de Dividendos}

\vspace{1.5cm}

{\color{azulul}\rule{0.6\linewidth}{0.8pt}}\\[0.5cm]
{\LARGE\textbf{Soluci\'on: Ejercicio 2}}\\[0.3cm]
{\Large\textit{Industrias LAN S.A.}}\\[0.3cm]
{\Large\textit{Dividendo vs.\ Recompra de Acciones (Irrelevancia MM)}}\\[0.5cm]
{\color{azulul}\rule{0.6\linewidth}{0.8pt}}

\vspace{2cm}

\begin{tcolorbox}[colback=azulclaro!30, colframe=azulul, width=0.75\linewidth, halign=center]
\centering
\textbf{Datos clave del problema}\\[0.4cm]
\begin{tabular}{ll}
Precio actual por acci\'on ($P_0$):      & \$20.00                \\
Acciones en circulaci\'on ($N_0$):       & 1{,}000                \\
Valor de mercado del equity ($E_0$):     & \$20{,}000             \\
Monto a distribuir:                      & \$1{,}000              \\
Dividendo por acci\'on (alternativa A):  & \$1.00                 \\
Recompra total (alternativa B):          & \$1{,}000              \\
Utilidad total anual:                    & \$2{,}000              \\
Impuestos:                               & Ignorar                \\
\end{tabular}
\end{tcolorbox}

\vfill
{\large Junio 2026}

{\color{azulul}\rule{\linewidth}{1.5pt}}
\end{titlepage}

% ========================
%  ENUNCIADO
% ========================
\begin{cajaenunciado}
\textbf{Industrias LAN S.A.} tiene el siguiente balance a valor de mercado.
Las acciones actualmente se venden a \textbf{\$20 la acci\'on}, y hay
\textbf{1{,}000 acciones en circulaci\'on}. La empresa pagar\'a un
\textbf{dividendo por acci\'on de \$1} o \textbf{recomprar\'a acciones
valoradas a \$1{,}000}. Ignore los impuestos.

\begin{enumerate}[label=\alph*.]
\item \textquestiondown Cu\'al ser\'a el \textbf{precio por acci\'on} bajo cada alternativa
      (dividendo vs.\ recompra)?
\item Si la \textbf{utilidad total de la empresa es \$2{,}000} al a\~no, encuentre la
      \textbf{Utilidad por Acci\'on (UPA)} bajo cada alternativa.
\end{enumerate}
\end{cajaenunciado}

\vspace{0.5cm}

% ========================
%  INTUICIÓN GENERAL
% ========================
\begin{cajaintuicion}{Dividendo vs.\ Recompra: dos caras de la misma moneda}
La \textbf{Proposici\'on de Irrelevancia de Dividendos
(Modigliani \& Miller, 1961)} establece que en mercados perfectos
---sin impuestos, sin costos de transacci\'on e informaci\'on
sim\'etrica--- la forma en que una empresa distribuye valor a sus
accionistas \emph{no altera} la riqueza total de estos.

En ambas alternativas, la empresa retira \textbf{\$1{,}000 de efectivo} del balance.
La diferencia es \emph{c\'omo} se entrega ese valor:

\begin{center}
\begin{tabular}{lll}
\toprule
 & \textbf{Dividendo} & \textbf{Recompra} \\
\midrule
Mecanismo   & Efectivo al accionista & Reducir acciones en circulaci\'on \\
Precio post & Baja en \$1 (ex-div.)  & Se mantiene en \$20 \\
Riqueza total & \$19 + \$1 = \$20    & \$20 sin cambio \\
\bottomrule
\end{tabular}
\end{center}
\end{cajaintuicion}

\vspace{1cm}

% ========================
%  SITUACIÓN INICIAL
% ========================
\section{Situaci\'on Inicial}

Antes de cualquier distribuci\'on, el balance a valor de mercado muestra:

\begin{align}
E_0 &= P_0 \times N_0 = \$20 \times 1{,}000 = \mathbf{\$20{,}000}
\end{align}

\noindent La empresa dispone de \textbf{\$1{,}000} para distribuir a sus accionistas,
ya sea como dividendo en efectivo (\textbf{Alternativa A}) o mediante la
recompra de acciones propias (\textbf{Alternativa B}).

\vspace{1cm}

% ========================
%  PARTE a — PRECIO
% ========================
\section{Parte a: Precio por Acci\'on bajo cada Alternativa}

% ---------- ALTERNATIVA A ----------
\subsection*{Alternativa A --- Pago de Dividendo (\$1 por acci\'on)}

\begin{cajaintuicion}{El precio cae exactamente en el valor del dividendo}
Al pagar el dividendo, el efectivo \textbf{sale de la empresa hacia los accionistas}.
El valor del equity disminuye en el monto distribuido. Por eso, en la fecha
ex-dividendo el precio baja en exactamente \$1: el accionista no gana ni pierde
riqueza, solo cambia la forma en que la tiene (efectivo vs.\ precio de la acci\'on).
\end{cajaintuicion}

\begin{cajaformula}
\[
P_{\text{ex-div}} = \frac{E_0 - \text{Dividendo total}}{N_0}
= \frac{E_0 - D \times N_0}{N_0}
\]
\end{cajaformula}

\begin{align}
\text{Dividendo total} &= \$1 \times 1{,}000 = \$1{,}000 \\[6pt]
E_{\text{post}} &= \$20{,}000 - \$1{,}000 = \$19{,}000 \\[6pt]
P_{\text{ex-div}} &= \frac{\$19{,}000}{1{,}000}
= \mathbf{\$19.00 \text{ por acci\'on}}
\end{align}

\noindent\textbf{Verificaci\'on de riqueza del accionista:}
\[
\text{Riqueza} = P_{\text{ex-div}} + D = \$19 + \$1 = \$20 \quad \checkmark
\]

\vspace{0.8cm}

% ---------- ALTERNATIVA B ----------
\subsection*{Alternativa B --- Recompra de Acciones (\$1{,}000)}

\begin{cajaintuicion}{La recompra reduce el n\'umero de acciones, no el precio}
Al recomprar acciones al precio de mercado, la empresa paga \$20 por cada acci\'on
que retira. El valor del equity baja en \$1{,}000, pero tambi\'en bajan las acciones
en circulaci\'on. Como ambos n\'umeros se reducen \emph{en la misma proporci\'on},
el precio por acci\'on \textbf{no cambia}. Los accionistas que no venden conservan
exactamente la misma riqueza que antes.
\end{cajaintuicion}

\begin{cajaformula}
\[
\Delta N = \frac{\text{Monto de recompra}}{P_0}
\qquad
P_{\text{post}} = \frac{E_0 - \text{Monto recompra}}{N_0 - \Delta N}
\]
\end{cajaformula}

\begin{align}
\Delta N &= \frac{\$1{,}000}{\$20} = 50 \text{ acciones recompradas} \\[6pt]
N_{\text{post}} &= 1{,}000 - 50 = 950 \text{ acciones} \\[6pt]
E_{\text{post}} &= \$20{,}000 - \$1{,}000 = \$19{,}000 \\[6pt]
P_{\text{post}} &= \frac{\$19{,}000}{950}
= \mathbf{\$20.00 \text{ por acci\'on}}
\end{align}

\noindent\textbf{Verificaci\'on de riqueza del accionista (que no vendi\'o):}
\[
\text{Riqueza} = P_{\text{post}} = \$20.00 \quad \checkmark
\]

\vspace{0.5cm}

% ---- Tabla comparativa de precios ----
\subsection*{Tabla comparativa --- Precio por acci\'on}

\begin{center}
\begin{tabular}{lrrr}
\toprule
\textbf{Concepto} & \textbf{Antes} & \textbf{Alt.\ A (Dividendo)} & \textbf{Alt.\ B (Recompra)} \\
\midrule
Valor equity           & \$20{,}000 & \$19{,}000 & \$19{,}000 \\
Acciones en circulaci\'on & 1{,}000  & 1{,}000    & 950        \\
\textbf{Precio por acci\'on} & \$20.00 & \$19.00  & \$20.00    \\
Riqueza por accionista & \$20.00    & \$19+\$1=\$20.00 & \$20.00 \\
\bottomrule
\end{tabular}
\end{center}

\[
\boxed{
\text{Alt.\ A: } P_{\text{ex-div}} = \$19.00
\qquad
\text{Alt.\ B: } P_{\text{post-recompra}} = \$20.00
}
\]

\vspace{1cm}

% ========================
%  PARTE b — UPA
% ========================
\section{Parte b: Utilidad por Acci\'on (UPA) bajo cada Alternativa}

\begin{cajaintuicion}{La recompra infla la UPA, pero no crea valor}
Tras la recompra hay \emph{menos acciones} entre las cuales repartir la misma
utilidad total (\$2{,}000). Por eso la UPA sube. Sin embargo, el \textbf{ratio
Precio/Utilidad (P/U)} es id\'entico en ambos casos: la mayor UPA se compensa
exactamente con el mayor precio, dejando el valor real del accionista inalterado.
\end{cajaintuicion}

\begin{cajaformula}
\[
\text{UPA} = \frac{\text{Utilidad total}}{N_{\text{acciones}}}
\qquad
\frac{P}{U} = \frac{P_{\text{acci\'on}}}{\text{UPA}}
\]
\end{cajaformula}

\subsection*{Alternativa A --- Dividendo}

\begin{align}
\text{UPA}_A &= \frac{\$2{,}000}{1{,}000} = \mathbf{\$2.00 \text{ por acci\'on}} \\[8pt]
\left(\frac{P}{U}\right)_A &= \frac{\$19.00}{\$2.00} = \mathbf{9.5\times}
\end{align}

\subsection*{Alternativa B --- Recompra}

\begin{align}
\text{UPA}_B &= \frac{\$2{,}000}{950} = \mathbf{\$2.1053 \text{ por acci\'on}} \\[8pt]
\left(\frac{P}{U}\right)_B &= \frac{\$20.00}{\$2.1053} = \mathbf{9.5\times}
\end{align}

\subsection*{Tabla comparativa --- UPA y ratio P/U}

\begin{center}
\begin{tabular}{lrr}
\toprule
\textbf{Concepto} & \textbf{Alt.\ A (Dividendo)} & \textbf{Alt.\ B (Recompra)} \\
\midrule
Utilidad total                    & \$2{,}000  & \$2{,}000  \\
Acciones en circulaci\'on         & 1{,}000    & 950        \\
\textbf{UPA}                      & \$2.00     & \$2.1053   \\
Precio por acci\'on               & \$19.00    & \$20.00    \\
\textbf{Ratio P/U}                & 9.5$\times$ & 9.5$\times$ \\
\bottomrule
\end{tabular}
\end{center}

\vspace{1cm}

% ========================
%  RESULTADO FINAL
% ========================
\begin{cajaresultado}
\begin{center}
\large
\begin{tabular}{p{6.5cm}p{4cm}p{4cm}}
\textbf{Concepto}                    & \textbf{Dividendo} & \textbf{Recompra} \\[4pt]
\hline\\[-6pt]
Precio por acci\'on post-distribuci\'on & \$19.00        & \$20.00 \\[4pt]
Acciones en circulaci\'on               & 1{,}000        & 950     \\[4pt]
Utilidad por Acci\'on (UPA)             & \$2.00         & \$2.11  \\[4pt]
Ratio Precio/Utilidad (P/U)             & 9.5$\times$    & 9.5$\times$ \\[4pt]
\textbf{Riqueza del accionista}         & \textbf{\$20}  & \textbf{\$20} \\
\end{tabular}
\end{center}
\end{cajaresultado}

\vspace{1cm}

% ========================
%  REFLEXIÓN
% ========================
\begin{cajareflexion}{\textquestiondown Por qu\'e la UPA sube con la recompra pero el valor no cambia?}
Es tentador concluir que la recompra es \emph{mejor} porque eleva la UPA de
\$2.00 a \$2.11. Sin embargo, esto es una ilusi\'on contable.

La UPA sube \textbf{solo porque hay menos acciones}, no porque la empresa
genere m\'as utilidades. El ratio P/U es \emph{id\'entico} en ambos casos:

\[
\frac{P}{U} = \frac{\$19.00}{\$2.00} = \frac{\$20.00}{\$2.1053} = 9.5\times
\]

Esto es la esencia de la \textbf{Proposici\'on de Irrelevancia de Dividendos
(Modigliani-Miller, 1961)}: en ausencia de imperfecciones de mercado, la
pol\'itica de distribuci\'on \emph{no crea ni destruye valor}. Lo que importa
es la capacidad de la empresa de generar utilidades, no la forma en que las
reparte.

\medskip
\textbf{Lecci\'on clave:} Una mayor UPA derivada de una recompra de acciones
no implica un mayor valor intr\'inseco. Los mercados descuentan esta mec\'anica
y el precio se ajusta de modo que la riqueza total del accionista permanece
constante (\$20 en ambas alternativas).
\end{cajareflexion}

\vspace{1cm}

% ========================
%  DIAGRAMA COMPARATIVO
% ========================
\newpage
\section*{Diagrama: Efecto sobre el Balance de Mercado}

\vspace{0.5cm}

\begin{center}
\resizebox{0.92\linewidth}{!}{%
\begin{tikzpicture}[
  box/.style={draw, rounded corners=4pt, minimum width=3.8cm,
              minimum height=1.2cm, align=center, font=\small},
  lbl/.style={font=\small\bfseries, align=center},
  arr/.style={-{Stealth}, thick}
]

% ---- Estado inicial ----
\node[box, fill=azulclaro!40, draw=azulul, minimum width=4.5cm] (ei)
  at (0,0) {Equity: \$20{,}000\\1{,}000 acc.\ $\times$ \$20};
\node[lbl, above=6pt] at (ei.north) {ANTES};

% ---- Flecha a Alt A ----
\node[box, fill=verdeclaro!60, draw=green!50!black, minimum width=4.5cm]
  (ea) at (-3.5,-3.5)
  {Equity: \$19{,}000\\1{,}000 acc.\ $\times$ \$19\\+ \$1 cash/acc.};
\node[lbl, above=6pt] at (ea.north) {ALTERNATIVA A\\Dividendo};

\draw[arr, blue!60!black] (ei.south) -- ++(0,-0.6) -| (ea.north)
  node[midway, above, font=\scriptsize\itshape] {$-\$1{,}000$ efectivo};

% ---- Flecha a Alt B ----
\node[box, fill=verdeclaro!60, draw=green!50!black, minimum width=4.5cm]
  (eb) at (3.5,-3.5)
  {Equity: \$19{,}000\\950 acc.\ $\times$ \$20\\(50 acc.\ recompradas)};
\node[lbl, above=6pt] at (eb.north) {ALTERNATIVA B\\Recompra};

\draw[arr, blue!60!black] (ei.south) -- ++(0,-0.6) -| (eb.north)
  node[midway, above, font=\scriptsize\itshape] {$-\$1{,}000$ efectivo};

% ---- Riqueza igual ----
\node[box, fill=amarilloclaro, draw=orange!70!black, minimum width=6cm,
      minimum height=1.4cm] (riq) at (0,-6.5)
  {\textbf{Riqueza del accionista: \$20 en ambos casos}\\
   \small A: \$19 precio + \$1 dividendo \quad B: \$20 precio};
\draw[arr, dashed, orange!80!black] (ea.south) -- ++(0,-0.4) -| (riq.west);
\draw[arr, dashed, orange!80!black] (eb.south) -- ++(0,-0.4) -| (riq.east);

\end{tikzpicture}}
\end{center}

\vspace{0.8cm}

% ========================
%  VERIFICACIÓN NUMÉRICA
% ========================
\section*{Verificaci\'on num\'erica compacta}

\begin{align*}
\textbf{Alt.\ A:}\quad
P_{\text{ex-div}} &= \frac{20{,}000 - 1{,}000}{1{,}000} = \$19.00
\quad;\quad
\text{UPA} = \frac{2{,}000}{1{,}000} = \$2.00
\quad;\quad
P/U = 9.5\times \\[8pt]
\textbf{Alt.\ B:}\quad
P_{\text{post}} &= \frac{20{,}000 - 1{,}000}{1{,}000 - \frac{1{,}000}{20}}
= \frac{19{,}000}{950} = \$20.00
\quad;\quad
\text{UPA} = \frac{2{,}000}{950} = \$2.11
\quad;\quad
P/U = 9.5\times
\end{align*}

\[
\Rightarrow\quad
\boxed{
\textbf{Precio: Alt.\ A = \$19 \ | \ Alt.\ B = \$20}
\qquad
\textbf{P/U: id\'entico en ambas (9.5}\times\textbf{)}
}
\]

\end{document}
